% GENERATED FILE. Edit canonical lessons, then run npm run generate:books.
% canonical-source-hash: fnv1a64:d7b9c3a193feeece
% canonical-lessons: TA-C179-payanulla, TA-C179-ullur, TA-C179-velinattu, TA-C179-telivana, TA-C179-unarcci

\chapter{Useful, Local, Foreign, Clear, Feeling}
\label{ch:ta-useful-local-foreign-clear-feeling}
\hlchaptermodality{179}

\begin{chapteropening}
\textbf{By the end of this chapter:} \emph{I can say useful, local, foreign, clear and a feeling.}

Complete the last lesson of chapter 179: I can say useful, local, foreign, clear and a feeling.
\end{chapteropening}

\section[\texorpdfstring{\ta{பயனுள்ள} (payaṉuḷḷa)}{பயனுள்ள (payaṉuḷḷa)}]{\ta{பயனுள்ள} (payaṉuḷḷa) — useful}

\label{lesson:TA-C179-payanulla}



\begin{quote}
Before the new one: say the Tamil for free of charge, then the Tamil for ripe.
\end{quote}

\subsection*{You'll want to know: \ta{பயனுள்ள}}
\textbf{\ta{பயனுள்ள}} — \emph{payaṉuḷḷa} — \textquotedblleft{}useful\textquotedblright{}.

\begin{grammarlens}[title={What you've built}]
One more word for this chapter.
\end{grammarlens}

\subsection*{Your turn}
\begin{itemize}
\raggedright
  \item \emph{Say it:} \emph{payaṉuḷḷa}
  \item \emph{Say it:} \emph{payaṉuḷḷa}, once more
  \item \emph{Say it:} say \emph{payaṉuḷḷa}
  \item \emph{Recall:} say the Tamil for free of charge, then the Tamil for ripe, then say \emph{payaṉuḷḷa} again
\end{itemize}

\begin{tcolorbox}[breakable,colback=teal!4,colframe=teal!35!black,title={Before you move on}]
What does \ta{பயனுள்ள} mean? (\textquotedblleft{}Useful\textquotedblright{}.) Say it once more.
\end{tcolorbox}
\section[\texorpdfstring{\ta{உள்ளூர்} (uḷḷūr)}{உள்ளூர் (uḷḷūr)}]{\ta{உள்ளூர்} (uḷḷūr) — local}

\label{lesson:TA-C179-ullur}



\begin{quote}
Before the new one: say the Tamil for ripe, then the Tamil for useful.
\end{quote}

\subsection*{You'll want to know: \ta{உள்ளூர்}}
\textbf{\ta{உள்ளூர்}} — \emph{uḷḷūr} — \textquotedblleft{}local\textquotedblright{}.

\begin{grammarlens}[title={What you've built}]
One more word for this chapter.
\end{grammarlens}

\subsection*{Your turn}
\begin{itemize}
\raggedright
  \item \emph{Say it:} \emph{uḷḷūr}
  \item \emph{Say it:} \emph{uḷḷūr}, once more
  \item \emph{Say it:} say \emph{uḷḷūr}
  \item \emph{Recall:} say the Tamil for ripe, then the Tamil for useful, then say \emph{uḷḷūr} again
\end{itemize}

\begin{tcolorbox}[breakable,colback=teal!4,colframe=teal!35!black,title={Before you move on}]
What does \ta{உள்ளூர்} mean? (\textquotedblleft{}Local\textquotedblright{}.) Say it once more.
\end{tcolorbox}
\section[\texorpdfstring{\ta{வெளிநாட்டு} (veḷināṭṭu)}{வெளிநாட்டு (veḷināṭṭu)}]{\ta{வெளிநாட்டு} (veḷināṭṭu) — foreign}

\label{lesson:TA-C179-velinattu}



\begin{quote}
Before the new one: say the Tamil for useful, then the Tamil for local.
\end{quote}

\subsection*{You'll want to know: \ta{வெளிநாட்டு}}
\textbf{\ta{வெளிநாட்டு}} — \emph{veḷināṭṭu} — \textquotedblleft{}foreign\textquotedblright{}.

\begin{grammarlens}[title={What you've built}]
One more word for this chapter.
\end{grammarlens}

\subsection*{Your turn}
\begin{itemize}
\raggedright
  \item \emph{Say it:} \emph{veḷināṭṭu}
  \item \emph{Say it:} \emph{veḷināṭṭu}, once more
  \item \emph{Say it:} say \emph{veḷināṭṭu}
  \item \emph{Recall:} say the Tamil for useful, then the Tamil for local, then say \emph{veḷināṭṭu} again
\end{itemize}

\begin{tcolorbox}[breakable,colback=teal!4,colframe=teal!35!black,title={Before you move on}]
What does \ta{வெளிநாட்டு} mean? (\textquotedblleft{}Foreign\textquotedblright{}.) Say it once more.
\end{tcolorbox}
\section[\texorpdfstring{\ta{தெளிவான} (teḷivāṉa)}{தெளிவான (teḷivāṉa)}]{\ta{தெளிவான} (teḷivāṉa) — clear}

\label{lesson:TA-C179-telivana}



\begin{quote}
Before the new one: say the Tamil for local, then the Tamil for foreign.
\end{quote}

\subsection*{You'll want to know: \ta{தெளிவான}}
\textbf{\ta{தெளிவான}} — \emph{teḷivāṉa} — \textquotedblleft{}clear\textquotedblright{}.

\begin{grammarlens}[title={What you've built}]
One more word for this chapter.
\end{grammarlens}

\subsection*{Your turn}
\begin{itemize}
\raggedright
  \item \emph{Say it:} \emph{teḷivāṉa}
  \item \emph{Say it:} \emph{teḷivāṉa}, once more
  \item \emph{Say it:} say \emph{teḷivāṉa}
  \item \emph{Recall:} say the Tamil for local, then the Tamil for foreign, then say \emph{teḷivāṉa} again
\end{itemize}

\begin{tcolorbox}[breakable,colback=teal!4,colframe=teal!35!black,title={Before you move on}]
What does \ta{தெளிவான} mean? (\textquotedblleft{}Clear\textquotedblright{}.) Say it once more.
\end{tcolorbox}
\section[\texorpdfstring{\ta{உணர்ச்சி} (uṇarcci)}{உணர்ச்சி (uṇarcci)}]{\ta{உணர்ச்சி} (uṇarcci) — a feeling, an emotion}

\label{lesson:TA-C179-unarcci}



\begin{quote}
Before the new one: say the Tamil for foreign, then the Tamil for clear.
\end{quote}

\subsection*{You'll want to know: \ta{உணர்ச்சி}}
\textbf{\ta{உணர்ச்சி}} — \emph{uṇarcci} — \textquotedblleft{}a feeling, an emotion\textquotedblright{}.

\begin{grammarlens}[title={What you've built}]
One more word for this chapter.
\end{grammarlens}

\subsection*{Your turn}
\begin{itemize}
\raggedright
  \item \emph{Say it:} \emph{uṇarcci}
  \item \emph{Say it:} \emph{uṇarcci}, once more
  \item \emph{Say it:} say \emph{uṇarcci}
  \item \emph{Recall:} say the Tamil for foreign, then the Tamil for clear, then say \emph{uṇarcci} again
\end{itemize}

\begin{tcolorbox}[breakable,colback=teal!4,colframe=teal!35!black,title={Before you move on}]
What does \ta{உணர்ச்சி} mean? (\textquotedblleft{}A feeling, an emotion\textquotedblright{}.) Say it once more.
\end{tcolorbox}
